\subsection{The distribution of a computational method}
\label{sec:r15methoddistribution}

A computational method includes its initialization, training draws, validation
rule, stopping rule, and response to failure. We specify these objects jointly.
Let $\mathcal S$ be the sixteen complete execution streams listed in the
protocol, fixed before the confirmation experiment, and let $\pi_{m,s}$ denote
the policy returned by method $m$ under stream $s$. The randomized implementation
draws $s$ uniformly from $\mathcal S$; the experiment executes every member of
this distribution. A failed fit returns the feasible analytical schedule and
remains in the denominator. Conditional on the stated software environment,
the stream specifies the complete training and checking computation. The
distribution does not include arbitrary unrecorded initializations or a new
draw of the stopping bank on re-execution.

Write $J(\pi)$ for the payoff in the original continuous-time economy,
averaged over the specified nine initial-state profiles. The method-level
objects are
\begin{align}
 G_m&=\frac{1}{|\mathcal S|}\sum_{s\in\mathcal S}
       \{J(\pi_{m,s})-J(\bar\pi)\},\label{eq:r15methodgain}\\
 \Delta_{m,b}&=\frac{1}{|\mathcal S|}\sum_{s\in\mathcal S}
       \{J(\pi_{m,s})-J(\pi_{b,s})\}.
       \label{eq:r15methodcontrast}
\end{align}
The first object is improvement on the schedule. The second is the direct
comparison of two randomized computational procedures with the same declared
stream distribution. Enumeration accounts for the entire training distribution;
there is no sampling error from an unobserved subset of these sixteen streams.
The remaining uncertainty concerns their independent payoff simulations and
the numerical transfer to continuous time. An assertion about this finite
distribution is not an assertion about all possible training initializations.

The stopping target is a gain of $0.0005$. At each of three declared work
checkpoints, fitting pauses for a new, independent certification calculation.
Training stops when its lower endpoint attains the target. Failed checks,
initialization, action search, checkpoint input and output, and the independent
final verification are included in the measured work. The gain $0.001$ is a
secondary final-attainment threshold; it is not a second measured stopping
experiment. A failure to attain either target is retained.

\subsection{Simultaneous inference for the finite method distribution}
\label{sec:r15methodinference}

Each returned policy receives $N$ fresh confirmation paths. The confirmation
bank is independent of fitting, validation, and every stopping decision.
Methods share initial profiles and innovations within a stream, whereas the
confirmation banks of distinct streams are independent. A direct contrast is
formed on each common path. Its common analytical schedule cancels, and its
transfer account contains the two policy allowances rather than two copies of
the schedule allowance, when both entries are simulated candidate gains.
An analytical fallback has the exact gain zero. A comparison with that
fallback retains the other candidate's full schedule-relative transfer
account; two analytical fallbacks have the exact difference zero.

For one registered endpoint, let $X_{sj}$ be the numerical path statistic,
$Y_{sj}=\max\{-C_s,\min(C_s,X_{sj})\}$, and let $\beta_s$ enclose the
difference between the continuous-time endpoint $\theta_s$ and
$\mathbb E Y_{sj}$. Thus $\beta_s$ includes the appropriate diffusion,
implementation, clipping-tail, and arithmetic accounts. Put
$C=\max_s C_s$, $n=|\mathcal S|N$, $\bar\beta=|\mathcal S|^{-1}\sum_s\beta_s$,
and denote the pooled mean and unbiased sample variance by $\bar Y$ and
$\widehat v$.

\begin{proposition}[Finite-distribution comparison]
\label{prop:r15methodinterval}
Conditional on the complete executed policy ensemble, suppose the
confirmation observations are independent across pairs $(s,j)$ and the stated
range and transfer accounts hold. For $0<\delta<1$, the interval
\begin{equation}
 \bar Y\ \pm\left\{
 \sqrt{\frac{2\widehat v\log(4/\delta)}{n}}
 +\frac{14C\log(4/\delta)}{3(n-1)}+\bar\beta\right\}
 \label{eq:r15methodinterval}
\end{equation}
covers $|\mathcal S|^{-1}\sum_s\theta_s$ with probability at least
$1-\delta$. The observations need not be identically distributed across
streams.
\end{proposition}

The confirmation family assigns $0.02$ across 238 two-sided events: in each
dimension, four schedule comparisons and three direct NBO comparisons, each
reported for sixteen streams and their uniform method mean. Online stopping
receives a separate $0.01$ across its 384 possible checks. The mechanism and
protected-observation families each receive $0.01$. Their union therefore has
error probability at most $0.05$ for the new experiment. The historical
experiment retains its separately stated confidence event. Unused checks do
not release probability for additional comparisons.

The practical-equivalence margin is fixed at $10^{-4}$ payoff units. For a
direct interval $[L,U]$, material superiority requires $L>10^{-4}$, and
practical equivalence requires $-10^{-4}<L\le U<10^{-4}$. Statistical
superiority, $L>0$, is reported separately. An interval containing zero does
not establish equivalence. Under log flow utility, the fixed margin corresponds
to a proportional consumption fee
$1-\exp\{-10^{-4}/A\}$, where
$A=(1-e^{-\rho T})/\rho$; for the stated primitives this is approximately
$1.0201$ basis points.

Simultaneous stream intervals also enclose the probability that the
randomized implementation attains a target $g$. If they are $[L_s,U_s]$,
the corresponding bounds are
\begin{equation}
 \frac{\#\{s:L_s\ge g\}}{|\mathcal S|}
 \ \le\ \frac{\#\{s:\theta_s\ge g\}}{|\mathcal S|}
 \ \le\ \frac{\#\{s:U_s\ge g\}}{|\mathcal S|}.
 \label{eq:r15attainment}
\end{equation}
No binomial sampling interval is attached to a population that has been
enumerated in full. Work is likewise reported over all streams, with
unattained runs retained in the distribution of work to target. A median over
successful runs alone is not the method's median.

The economic decomposition records production, consumption, and terminal
dispersion contributions on the same numerical confirmation paths. These
components add to the observed numerical payoff gain and help describe the
policy's behavior. They have no separate continuous-time confidence statement:
the transfer allowance is attached to the total payoff and is not divided
among components without additional bounds.
